\section{Component Analysis of Two-Phase Hybrids}
\label{sec:ablation}
%% generated by mpce_results.py -- do not edit by hand
\begin{table}[!t]
\centering
\caption{Ablation over the 68 Benchmark Cases (30 Seed-Paired Runs, 6,030 Calls). Top: Average Rank Among the Eight Variants (Friedman $\chi^2_F=326.8$, 7 d.f., $p=1.14\times10^{-66}$) and Feasible Runs. Bottom: Planned Contrasts, Cases in Which the First Variant Is Significantly Better / Not Different / Worse (Wilcoxon, Holm-Adjusted over the Eight Contrasts of Each Case, $\alpha=0.05$), and Mean Difference of the Wake Loss $\overline{\Delta L}$ (Percentage Points; Negative = First Variant Better)}
\label{tab:ablation}
\scriptsize\setlength{\tabcolsep}{2.5pt}
\begin{tabular}{lcccc}
\toprule
Variant & Phase 1 & Phase 2 & Avg.\ rank & Feas.\ (\%) \\
\midrule
PSO-VNS & PSO & VNS & 1.87 & 100.0 \\
SSA-VNS & SSA & VNS & 3.79 & 99.6 \\
LX-SSA-VNS & LX-SSA & VNS & 4.62 & 99.9 \\
RS-VNS & random sampling & VNS & 4.46 & 99.4 \\
VNS & -- & VNS & 5.32 & 100.0 \\
PSO & PSO & -- & 2.10 & 100.0 \\
SSA & SSA & -- & 6.56 & 97.6 \\
LX-SSA & LX-SSA & -- & 7.29 & 96.9 \\
\bottomrule
\end{tabular}

\smallskip
\begin{tabular}{l>{\raggedright\arraybackslash}p{2.9cm}cc}
\toprule
Contrast & Isolates & W/T/L & $\overline{\Delta L}$ \\
\midrule
PSO-VNS vs.\ PSO & VNS phase & 8/58/2 & $-0.018$ \\
PSO-VNS vs.\ VNS & swarm start vs.\ best initial point & 50/17/1 & $-0.465$ \\
PSO-VNS vs.\ RS-VNS & swarm vs.\ random sampling & 48/20/0 & $-0.399$ \\
PSO-VNS vs.\ SSA-VNS & PSO vs.\ SSA as Phase~1 & 47/21/0 & $-0.330$ \\
PSO-VNS vs.\ LX-SSA-VNS & PSO vs.\ LX-SSA as Phase~1 & 48/20/0 & $-0.393$ \\
SSA-VNS vs.\ RS-VNS & SSA vs.\ random sampling & 2/66/0 & $-0.068$ \\
SSA-VNS vs.\ LX-SSA-VNS & Laplace step in the hybrid & 3/65/0 & $-0.063$ \\
LX-SSA vs.\ SSA & Laplace step alone & 0/61/7 & $+0.166$ \\
\bottomrule
\end{tabular}
\end{table}


The eight variants of Table~\ref{tab:ablation} combine a Phase~1 (PSO, SSA, LX-SSA, random sampling or none) with or without the same VNS phase, at 6,030 calls and with the same seeds (convergence in Fig.~\ref{fig:ablation}). The Friedman test rejects equal average ranks ($\chi^2_F=\NAblChi$, 7 d.f., $p=\NAblP$).

\emph{What the VNS phase contributes.} Adding VNS helps every swarm: each hybrid ranks ahead of the same swarm run alone for the full budget (PSO-VNS vs.\ PSO, SSA-VNS vs.\ SSA and LX-SSA-VNS vs.\ LX-SSA in Table~\ref{tab:ablation}); e.g., the mean wake loss falls from \NLossSSA\% (SSA) to \NLossSSAVNS\% (SSA-VNS). PSO-VNS and PSO share their first 3,030 calls, so their contrast measures only the use of the second half: switching to VNS instead of continuing the swarm gives a small gain that grows with $N$ (\NAblPSOVNSvsPSO{} under the Holm adjustment over the contrasts of Table~\ref{tab:ablation}, mean difference \NAblPSOVNSvsPSODL{} percentage points; Section~\ref{sec:overall}): a convergent PSO keeps improving in the second half on small layouts, whereas on larger layouts the single-turbine moves of the compass search pay off.
% CHECK-FINAL [C13]: ablation.contrasts.PSOBV-PSOC W > L

\emph{What the swarm phase contributes.} Two controls isolate Phase~1. VNS alone starts from the best initial layout and spends the whole budget on refinement; RS-VNS spends the 3,015 calls of Phase~1 on independent random layouts and refines the best one. PSO-VNS is significantly better than VNS in \NAblPSOVNSvsVNSW{} cases and worse in \NAblPSOVNSvsVNSL{} (\NAblPSOVNSvsVNSDL{} percentage points), and better than RS-VNS in \NAblPSOVNSvsRSVNSW{} cases and never worse (\NAblPSOVNSvsRSVNSDL{} percentage points). The PSO phase thus adds value beyond a longer refinement and beyond a random multistart of equal cost. The salp-swarm phases do not: SSA-VNS and RS-VNS are statistically tied in \NAblSSAVNSvsRSVNST{} of the 68 cases (\NAblSSAVNSvsRSVNS, \NAblSSAVNSvsRSVNSDL{} percentage points), and LX-SSA-VNS ranks even behind RS-VNS on average. As start generators for VNS, SSA and LX-SSA are no better than random sampling, which itself ranks ahead of VNS alone.
% CHECK-FINAL [C14]: PSOBV-BVNS W >= 40, L <= 2; PSOBV-RSVNS W >= 40, L == 0
% CHECK-FINAL [C15]: SSABV-RSVNS T >= 60, W and L <= 5

\emph{Laplace step.} The Laplace step of LX-SSA gives no gain: without the VNS phase, LX-SSA is \NAblLXSSAvsSSA{} against SSA, and inside the hybrid, SSA-VNS is \NAblSSAVNSvsLXSSAVNS{} against LX-SSA-VNS (\NAblSSAVNSvsLXSSAVNSDL{} percentage points).
% CHECK-FINAL [C17]: LXSSA-SSA W == 0; SSABV-LXBV L == 0

\emph{Choice of Phase~1.} Among the Phase-1 options tested, only PSO is effective. With the same VNS phase, PSO-VNS is \NAblPSOVNSvsSSAVNS{} (W/T/L) against SSA-VNS and \NAblPSOVNSvsLXSSAVNS{} against LX-SSA-VNS. At the switch, its global best is feasible more often (\NSwitchFeasPSOVNS\% of the runs vs.\ \NSwitchFeasSSAVNS\% and \NSwitchFeasLXSSAVNS\%) and, when feasible, has a lower wake loss (\NSwitchLossPSOVNS\% vs.\ \NSwitchLossSSAVNS\% and \NSwitchLossLXSSAVNS\%), so VNS starts from a better layout. The first phase of a hybrid should therefore be judged against a random-sampling control of equal cost.
% CHECK-FINAL [C16]: PSOBV-SSABV, PSOBV-LXBV: L == 0, W >= 34; switch feasibility higher and switch loss lower for PSOBV

\subsection{Budget Split}
\label{sec:split}
%% generated by mpce_results.py -- do not edit by hand
\begin{table}[!t]
\centering
\caption{Sensitivity of PSO-VNS to the budget split between the PSO and VNS phases (25\%, 50\% and 75\% of the 6,030 calls for PSO): mean benchmark objective of the feasible runs (superscript: feasible runs when fewer than 30) and Holm-adjusted Wilcoxon signed-rank $p$ of the 50\% split against the 25\% and 75\% splits (30 seed-paired runs).}
\label{tab:split}
\scriptsize\setlength{\tabcolsep}{2.5pt}
\begin{tabular}{cccccccc}
\toprule
DS & $r$ & $N$ & 25\% & 50\% & 75\% & $p$ (25) & $p$ (75) \\
\midrule
I & 500 & 6 & 83851.3 & 83918.5 & 83929.5 & 0.023 & 0.465 \\
I & 500 & 10 & 135131.6 & 135160.5 & 135421.8 & 0.556 & 0.069 \\
I & 750 & 8 & 111835.7 & 111921.0 & 111944.0 & 0.014 & 0.045 \\
I & 750 & 12 & 164894.0 & 165526.4 & 165991.2 & 0.004 & 0.004 \\
I & 1000 & 10 & 139700.6 & 139874.2 & 139883.7 & 0.002 & 0.529 \\
I & 1000 & 15 & 207367.9 & 207308.1 & 207724.0 & 0.570 & $3.20\times10^{-5}$ \\
II & 500 & 6 & 42809.4 & 42851.5 & 42883.5 & 0.110 & 0.110 \\
II & 500 & 10 & 67624.1 & 67773.5 & 67799.8 & 0.115 & 0.903 \\
II & 750 & 8 & 57358.0 & 57411.1 & 57448.4 & 0.121 & 0.121 \\
II & 750 & 12 & 82784.6 & 82968.9 & 83182.0 & 0.070 & 0.006 \\
II & 1000 & 10 & 71769.3 & 71906.7 & 71976.0 & 0.015 & 0.013 \\
II & 1000 & 15 & 104349.6 & 104794.8 & 104784.9 & $4.63\times10^{-4}$ & 0.968 \\
\bottomrule
\end{tabular}
\end{table}

The split $\omega=0.5$ was fixed before the experiments and is used for all other results; Table~\ref{tab:split} compares it with $\omega=0.25$ and $0.75$ on the \NSplitCases{} cases with the middle and the largest $N$ of each farm (per-case values in Table~\ref{S-tab:split-cases} of the supplementary material). A longer PSO phase is better: with $\omega=0.75$, the mean wake loss is lower than with $\omega=0.5$ in \NSplitSeventyFiveLower{} of the \NSplitCases{} cases, significantly in \NSplitSeventyFiveSig, and \NSplitGainSeventyFive{} percentage points lower on average, and the default is significantly better in \NSplitFiftyBetterSeventyFive{} case; $\omega=0.25$ is significantly worse than the default in \NSplitTwentyFiveSig{} cases and better in \NSplitTwentyFiveBetter{} case. In line with the component analysis, most of the benefit comes from the PSO phase, and VNS serves as a final refinement. We do not re-tune $\omega$ on the test cases, so the main comparison may slightly understate PSO-VNS; since the split study covers only \NSplitCases{} cases and three values, we give $\omega=0.75$ as a recommended starting point for applications, not as an optimized setting.
% CHECK-FINAL [C18]: split: 75% lower loss in >= 10 of 12; 50% never significantly better than 75%; 25% never significantly better than 50%; ranks 75 < 50 < 25
